\documentclass[11pt]{article}
% Source-PDF: 数学 MATHEMATICS/快速傅里叶变换学习总结 - Voiphon.pdf
\usepackage{oipl-vn}

\title{Tổng kết học nhanh về biến đổi Fourier nhanh}
\author{Voiphon}
\date{}

\begin{document}
\maketitle

\origtitle{Fast Fourier Transform study summary}
\sourcepath{Mathematics / Fast Fourier Transform study summary - Voiphon.pdf}

\section{Mở đầu}

Tác giả đã từng đọc nội dung này trong \textit{Introduction to Algorithms} từ
lâu, nhưng lúc đó chưa thật sự hiểu. Khi tìm lại tài liệu trên mạng, tác giả
thấy không có nhiều bài viết đủ đơn giản và phù hợp với người mới học; phần
lớn tài liệu nghiêng về phân tích tín hiệu số. Vì vậy bài viết này được ghi
lại như một bản tổng kết để tránh quên kiến thức sau vài ngày.

\section{Tích chập: phép nhân đa thức}

Định nghĩa một đa thức bậc $n$ theo dạng
\[
  f(x)=a_0+a_1x+\cdots+a_{n-1}x^{n-1}.
\]
Ở đây cách gọi "bậc $n$" hơi không chuẩn, vì chỉ số cao nhất là $n-1$.

Tích của hai đa thức $f(x)$ và $g(x)$ hiển nhiên vẫn là một đa thức. Ta gọi
tích này là \term{tích chập} của $f(x)$ và $g(x)$. Nghiêm túc hơn, tích chập
thật sự thường được định nghĩa bằng tích phân trên các hàm liên tục, nhưng
trong bối cảnh lập trình thi đấu, điều cần tính thường là tích chập của hai
dãy hệ số.

Cách tính trực tiếp là nhân đa thức:
\[
  h(x)=f(x)g(x)
  =a_0b_0+(a_0b_1+a_1b_0)x+\cdots+
  \left(\sum_{i=0}^{k} a_i b_{k-i}\right)x^k+\cdots.
\]
Tách hệ số của $x^k$, ta có
\[
  c_k=\sum_{i=0}^{k} a_i b_{k-i}.
\]
Đó chính là tích chập của hai dãy số. Cách tính thô trực tiếp theo công thức
trên có độ phức tạp $O(n^2)$.

\section{Biểu diễn điểm, tính giá trị và nội suy đa thức}

Một đa thức với $n$ hệ số không chỉ có thể biểu diễn bằng bộ hệ số
\[
  (a_0,a_1,\ldots,a_{n-1}),
\]
mà còn có thể biểu diễn bằng $n$ điểm phân biệt
\[
  (x_0,y_0),(x_1,y_1),\ldots,(x_{n-1},y_{n-1}).
\]
Lý do là đa thức này có đúng $n$ hằng số cần xác định.

Ưu điểm của biểu diễn điểm là tính tích chập rất nhanh: với những điểm có cùng
hoành độ $x$, ta chỉ cần nhân trực tiếp các giá trị hàm tương ứng. Độ phức tạp
khi đó là $O(n)$.

Quá trình chuyển từ biểu diễn hệ số sang biểu diễn điểm được gọi là
\term{tính giá trị} (evaluation). Cách trực tiếp là thay lần lượt $n$ điểm vào
đa thức. Với mỗi điểm có thể tính theo dạng Horner:
\[
  f(x)=a_0+x(a_1+x(a_2+\cdots)).
\]
Mỗi điểm tốn $O(n)$, tổng cộng $O(n^2)$.

Quá trình chuyển từ biểu diễn điểm về biểu diễn hệ số được gọi là
\term{nội suy} (interpolation). Có thể làm bằng phương pháp hệ số chưa biết
hoặc nội suy Lagrange; bài viết này không dùng trực tiếp phần đó.

\section{\texorpdfstring{DFT: biến đổi Fourier rời rạc trong $O(n^2)$}{DFT: biến đổi Fourier rời rạc trong O(n\textasciicircum 2)}}

Các điểm $x$ của đa thức không nhất thiết phải là số thực. DFT, viết tắt của
\term{Discrete Fourier Transform}, thực chất là việc tính giá trị của một đa
thức tại các căn bậc $n$ của đơn vị.

Căn bậc $n$ của đơn vị là $n$ số phức
\[
  \omega_n^k=e^{2\pi ki/n}
  =\cos\left(k\cdot\frac{2\pi}{n}\right)
  + i\sin\left(k\cdot\frac{2\pi}{n}\right).
\]
Nói trực quan, đó là tọa độ của $n$ đỉnh của đa giác đều nội tiếp đường tròn
đơn vị, với một đỉnh tại $(1,0)$.

Công thức DFT là
\[
  \operatorname{DFT}(k)=y_k=f(\omega_n^k)
  =\sum_{i=0}^{n-1} a_i(\omega_n^k)^i.
\]
Cách tính trực tiếp theo công thức này có độ phức tạp $O(n^2)$.

\section{\texorpdfstring{FFT: biến đổi Fourier nhanh trong $O(n\log n)$}{FFT: biến đổi Fourier nhanh trong O(n log n)}}

FFT là thuật toán tính DFT trong thời gian $O(n\log n)$. Ý tưởng chính là khai
thác tính chất đặc biệt của căn bậc $n$ của đơn vị, rồi dùng chia để trị để
tăng tốc.

Giả sử $n$ là số chẵn. Khi đó
\[
  (\omega_n^k)^2=(\omega_n^{k+n/2})^2=\omega_{n/2}^{k}.
\]
Xét đa thức
\[
  f(x)=a_0+a_1x+a_2x^2+\cdots+a_{n-1}x^{n-1}.
\]
Ta tách nó thành phần hệ số chẵn và phần hệ số lẻ:
\[
  f_1(x)=a_0+a_2x^2+a_4(x^2)^2+\cdots+a_{n-2}(x^2)^{n/2-1},
\]
\[
  f_2(x)=a_1+a_3x^2+a_5(x^2)^2+\cdots+a_{n-1}(x^2)^{n/2-1}.
\]
Rõ ràng
\[
  f(x)=f_1(x)+x f_2(x).
\]

Điểm mấu chốt là khi tính trên các căn bậc $n$ của đơn vị, các giá trị $x^2$
chỉ còn $n/2$ giá trị khác nhau. Vì vậy $f_1$ và $f_2$ tổng cộng chỉ cần tính
trên $n$ điểm, thay vì $2n$ điểm như trường hợp các $x$ tùy ý. Nếu $n$ là lũy
thừa của $2$, ta tiếp tục chia để trị như vậy và đạt độ phức tạp
$O(n\log n)$. Nếu $n$ chưa phải lũy thừa của $2$, ta có thể thêm các hệ số
$0$ vào đa thức.

\section{IDFT: biến đổi Fourier rời rạc ngược}

DFT dùng để tính giá trị. Biến đổi ngược IDFT dùng để nội suy, tức chuyển từ
biểu diễn điểm về biểu diễn hệ số.

Viết DFT ở dạng ma trận, ta có một ma trận Vandermonde nhân với vector hệ số
bằng vector giá trị:
\[
  V_n
  \begin{pmatrix}
    a_0\\a_1\\a_2\\ \vdots \\ a_{n-1}
  \end{pmatrix}
  =
  \begin{pmatrix}
    y_0\\y_1\\y_2\\ \vdots \\ y_{n-1}
  \end{pmatrix}.
\]
Ma trận bên trái là ma trận Vandermonde. Nghịch đảo của nó đã có công thức
chuẩn, từ đó suy ra công thức IDFT:
\[
  \operatorname{IDFT}(k)=a_k
  =\frac{1}{n}\sum_{i=0}^{n-1} y_i(\omega_n^{-k})^i.
\]

So sánh với công thức DFT
\[
  \operatorname{DFT}(k)=y_k=\sum_{i=0}^{n-1} a_i(\omega_n^k)^i,
\]
ta thấy hai công thức có hình thức rất giống nhau. Vì vậy chỉ cần thay tham số
trong FFT, ta cũng có thể nội suy trong $O(n\log n)$.

\section{Quay lại tích chập}

Ta có thể tính tích chập bằng công thức
\[
  f\cdot g = \operatorname{IDFT}(\operatorname{DFT}(f)\cdot
  \operatorname{DFT}(g)).
\]
Nghĩa là trước hết tính giá trị để chuyển hai đa thức sang biểu diễn điểm, sau
đó nhân theo từng điểm, rồi nội suy về đa thức mới.

Khi dùng FFT, tổng thời gian là
\[
  O(n\log n)+O(n)+O(n\log n)=O(n\log n).
\]
Đây chính là thuật toán nhanh cho tích chập: FFT.

Một mở rộng thường gặp là NTT (Number Theoretic Transform, biến đổi số học),
cũng có độ phức tạp tương tự nhưng tránh sai số số thực và thường có hằng số
tốt hơn trong một số bài toán.

\section{Ứng dụng}

Ví dụ mẫu: CodeVS 3123 Super A*B Problem. Cho hai số $A$ và $B$, cần tính
$A\cdot B$, trong đó $A$ và $B$ có thể dài tới $10^5$ chữ số.

Phân tích: một số nguyên lớn có thể xem như một đa thức tại $x=10$, hoặc trực
tiếp xem các chữ số của hai số là hai dãy cần lấy tích chập.

\section{Mẫu mã}

PDF gốc có các mục "mẫu mã đệ quy" và "mẫu mã không đệ quy", nhưng phần mã
không được trích xuất rõ bằng \texttt{pdftotext}. Vì vậy bản dịch này giữ lại
nội dung lý thuyết và ghi chú rằng phần mã cần được phục dựng riêng nếu có
nguồn gốc tốt hơn PDF hiện tại.

\section*{Ghi chú nguồn}

PDF gốc dẫn hai bài viết CSDN về FFT và ghi chú rằng sự giống nhau giữa DFT và
IDFT không phải ngẫu nhiên: chúng là hai quy tắc biến đổi đối xứng như ảnh
gương trên miền hàm phức.

\end{document}
